\documentclass{article}
\usepackage[UTF8]{ctex}
\usepackage{amsmath,amssymb,amsfonts}
\usepackage{enumitem}
\usepackage{geometry}
\geometry{a4paper, margin=1in}
% === 错题本排版优化（不参与题目标识） ===
\setlength{\parindent}{0pt}
\setlength{\parskip}{0.35em plus 0.10em minus 0.08em}
\setlength{\emergencystretch}{3em}
\setlist{itemsep=0.25em, topsep=0.40em, parsep=0pt}
\allowdisplaybreaks[2]
% === 错题本排版优化结束 ===
\begin{document}
\section*{西北大学2025年高等代数真题}
\section*{167．西北大学2025年高等代数试卷}

\subsection{题目1}
1．（10 分）已知 \(n\) 为正整数，求多项式 \(x^{2}-x+1\) 与 \(x^{n+2}-(x-1)^{2 n+1}\) 的首一最大公因式．

\subsection{题目2}
2．（15 分）\(a, b, c\) 均为整数，\(a c+b c\) 为奇数，将多项式

\[
x^{6}-(6-a) x^{5}+(12-6 a+b) x^{4}+(12 a-8-6 b+c) x^{3}+(-8 a+12 b-6 c) x^{2}+(12 c-8 b) x-8 c
\]

分解为有理数域上不可约多项式的乘积。

\subsection{题目3}
3．（10 分）计算行列式

\[
f(x)=\left|\begin{array}{cccccc}
x & x & \cdots & x & x & x \\
x & x & \cdots & x & a_{1} & x \\
2 x & 2 x & \cdots & a_{2} & 2 x & 2 x \\
\vdots & \vdots & & \vdots & \vdots & \vdots \\
a_{n} & n x & \cdots & n x & n x & n x
\end{array}\right|
\]

并求 \(f(x)\) 的根．

\subsection{题目4}
4．（15 分）已知 \(\alpha_{1}, \alpha_{2}, \cdots, \alpha_{n}\) 为 \(n\) 维线性空间 \(V\) 的一组线性无关的向量，讨论 \(\beta_{1}, \beta_{2}, \cdots, \beta_{n}\) 的线性相关性，其中 \(\beta_{i}=(n+1-i) \alpha_{i}+i \alpha_{i+1}, i=1,2, \cdots, n-1, \beta_{n}=\alpha_{n}+n \alpha_{1}\) ．

\subsection{题目5}
5．（10 分）已知 \(n+1(n \geq 2)\) 阶方阵 \(A\) 的逆矩阵为

\[
A^{-1}=\left(\begin{array}{ccccc}
1 & a & a^{2} & \cdots & a^{n} \\
0 & 1 & a & \cdots & a^{n-1} \\
0 & 0 & 1 & \cdots & a^{n-2} \\
\vdots & \vdots & \vdots & & \vdots \\
0 & 0 & 0 & \cdots & 1
\end{array}\right)
\]

\(A^{*}\) 为 \(A\) 的伴随矩阵，求 \(\left(A^{*}\right)^{*}\) 。

\subsection{题目6}
6．（15 分）当 \(a, b\) 为何值时，线性方程组

\[
\left\{\begin{array}{l}
x_{1}+x_{2}+x_{3}+x_{4}+x_{5}=1 \\
3 x_{1}+2 x_{2}+x_{3}+x_{4}-3 x_{5}=a \\
x_{2}+2 x_{3}+2 x_{4}+6 x_{5}=3 \\
5 x_{1}+4 x_{2}+3 x_{3}+3 x_{4}-x_{5}=b
\end{array}\right.
\]

无解？有解？并在有解时求解。

\subsection{题目7}
7．（15 分）已知 \(A, B\) 均为 \(n\) 阶正定矩阵，证明：

\begin{itemize}
  \item[（1）] \(A+2025 B\) 为正定矩阵。
  \item[（2）] \(|A+2025 B|>|A|\) ．
\end{itemize}

\subsection{题目8}
8．（15 分）已知 \(n\) 维线性空间 \(V\) 中 \(\alpha_{1}, \alpha_{2}, \cdots, \alpha_{r}\) 线性无关，\(\beta_{1}, \beta_{2}, \cdots, \beta_{n}\) 线性无关，且 \(r+s>n\) 。证明：存在 \(\xi \in V\) ，使得 \(\xi\) 既能由 \(\alpha_{1}, \alpha_{2}, \cdots, \alpha_{r}\) 线性表出，也能由 \(\beta_{1}, \beta_{2}, \cdots, \beta_{s}\) 线性表出。

\subsection{题目9}
9．（15 分）已知 \(M_{n}(F)\) 为数域 \(F\) 上的 \(n\) 阶矩阵构成的线性空间，\(A, B \in M_{n}(F)\) ．定义

\[
\varphi: \varphi(X)=A X B, X \in M_{n}(F) .
\]

\begin{itemize}
  \item[（1）] 证明：\(\varphi\) 是线性变换。
  \item[（2）] 证明：\(\varphi\) 可逆的充要条件是 \(A, B\) 为可逆矩阵．
  \item[（3）] 若 \(r(A)=r, r(B)=s\) ，求 \(\operatorname{Ker} \varphi\) 的维数和基．
\end{itemize}

\subsection{题目10}
10．（15 分）\(A\) 为 \(n\) 阶方阵，存在正整数 \(m\) 使得 \(A^{m-1} \neq O, A^{m}=O\) ，则 \(A\) 为 \(n\) 阶 \(m\) 次幂零矩阵。证明：所有的 \(n\) 阶 \(n-1\) 次幂零矩阵都相似，并求其若尔当标准形。

\subsection{题目11}
11．（15 分）二次型

\[
f\left(x_{1}, x_{2}, x_{3}\right)=x_{1}^{2}+2 x_{2}^{2}+3 x_{3}^{2}-4 x_{1} x_{2}+2 a x_{2} x_{3}(a>0)
\]

在正交变换 \(X=Q Y\) 下的标准形为 \(f=2 y_{1}^{2}-y_{2}^{2}+5 y_{3}^{2}\) ，求参数 \(a\) 及相应的正交矩阵 \(Q\) 。
\end{document}
